Autoregressive language models train on all tokens in parallel, yet must generate them one at a time.
Diffusion language models remove this limit by denoising many tokens in each forward pass \citep{nie2026large, ye2025dream7bdiffusionlarge, team2026diffusiongemma}.
Samplers are a natural place to improve these models: in continuous diffusion, many advances have come from better samplers \citep{song2020denoising}, and in language diffusion, the sampling order strongly affects task performance \citep{kim2025train, ben2026accelerated, kim2026klass}.

Any diffusion language sampler makes two choices for each position in its denoising window.
The first is \emph{where} the next token comes from: the model's prediction, the token already at that position, or fresh noise.
This defines when a position within a denoising window has converged.
The second is \emph{when} the position is committed to the KV cache, after which it can no longer change.
These design choices give rise to two potential extremes: on one end, it is possible to commit tokens on a \textit{per position} basis, or as soon as a position has converged.
It is also possible to commit tokens \textit{at a block level}, or after the entire denoising window has converged.

In this work, we use the Entropy Bound sampler used in DiffusionGemma \citep{team2026diffusiongemma} to study these extremes and uncover the underlying tradeoffs.
Committing each token as soon as it converges removes the chance to revise it later.
Revision matters most when the entropy bound is relaxed to trade accuracy for speed.
When accepted tokens are frozen, macro accuracy falls from 73\% to 20\% as the entropy bound is relaxed, while the unfrozen sampler stays near 70\%.
Committing on a block level results in many redundant computations, as tokens converge in a largely left-to-right order.
Thus a converged prefix sits idle in the window, and by twenty denoising steps more than half of the slots, on average, are spent on tokens that are already final.

Since per-position commitment gives up revision and block-level commitment wastes computation, we introduce the \textbf{R}olling \textbf{E}ntropy \textbf{S}ampler (RES).
RES commits a contiguous prefix of accepted positions once that prefix holds at least $\ell_{\min}$ positions, and it keeps positions farther ahead open to revision.
Earlier rolling samplers work with masked tokens, which are frozen once accepted.
DSB \citep{luo2026dsb} rolls the window for masked diffusion, and Set Block Decoding \citep{gat2025set} fine-tunes an autoregressive model to decode masked tokens ahead.
Concurrent work on Clock Diffusion \citep{schiff2026clock} instead trains a Gaussian diffusion model with a noise schedule that depends on position.
RES differs from both in two ways.
First, it runs on uniform diffusion, where every position always holds a token that can still change. Section~\ref{sec:too-early} shows that keeping this ability matters.
Second, RES changes only inference, so it needs no additional training.

In summary, our contributions are as follows:
\begin{itemize}
    \item Separating accepting a token from committing it (Section~\ref{sec:motivation}), we show that DiffusionGemma revises its own accepted tokens without renoising, and that freezing them lowers macro accuracy from 73\% to 20\% at loose entropy bounds (Section~\ref{sec:too-early}).
    \item We show that the EB sampler converges left to right, leaving over half of the window idle late in each block (Section~\ref{sec:too-late}).
    \item We introduce RES, a rolling sampler that needs no training (Section~\ref{sec:method}). Across seven benchmarks it increases tokens per forward pass from 19.3 to 25.5, while macro accuracy goes from 80.7\% to 78.8\%. It also extends the speed-quality frontier beyond what relaxing the entropy bound alone reaches (Section~\ref{sec:experiments}).
\end{itemize}
